\section*{Chapter \ref{ch:g12:exp} --- Exponential and Logarithm}

\begin{solution}{exo:g12:exp:1}
$\dfrac{\eu^{3x}\,\eu^{-x+1}}{\eu^{x}} = \eu^{3x - x + 1 - x} = \eu^{x+1}$.

$\ln\!\left(\dfrac{\eu^2\sqrt{\eu}}{\eu^{-3}}\right)
= 2 + \tfrac12 + 3 = \tfrac{11}{2}$.
\end{solution}

\begin{solution}{exo:g12:exp:2}
\emph{(a)} Set $X = \eu^x > 0$: $X^2 - 3X + 2 = 0$ gives $X = 1$ or $X = 2$,
both positive, so $x = 0$ or $x = \ln 2$.

\emph{(b)} \omterm{def:g11:func:function}{Domain}: $x - 1 > 0$ and $x + 2 > 0$, so $x > 1$. The \omterm{def:g10:algebra:equation}{equation}
becomes $\ln\bigl((x-1)(x+2)\bigr) = \ln 4$, hence $(x-1)(x+2) = 4$,
\emph{i.e.} $x^2 + x - 6 = 0$, so $x = 2$ or $x = -3$. Only $x = 2$ is in the
\omterm{def:g11:func:function}{domain}: $S = \{2\}$.
\end{solution}

\begin{solution}{exo:g12:exp:3}
Divide by $\eu^x$:
$\dfrac{1 - x^2\eu^{-x}}{1 + \eu^{-x}} \to \dfrac{1-0}{1+0} = 1$, using
$x^2 \eu^{-x} \to 0$ (\cref{thm:g12:exp:growth}).

$\dfrac{\ln x}{\sqrt x} \to 0$ by \cref{prop:g12:exp:lnanalytic} (case
$n = 2$).

$x^2 \ln x = x \cdot (x \ln x) \to 0 \times 0 = 0$.

$x - \ln x = x\left(1 - \frac{\ln x}{x}\right) \to +\infty$ since
$\frac{\ln x}{x} \to 0$.
\end{solution}

\begin{solution}{exo:g12:exp:4}
$f'(x) = \eu^{-x} - x\eu^{-x} = (1 - x)\eu^{-x}$, of the sign of $1 - x$:
$f$ increases on $\intoc{-\infty}{1}$, decreases on $\intco{1}{+\infty}$,
with a global \omterm{def:g10:functions:extrema}{maximum} $f(1) = \eu^{-1}$.

Limits: as $x \to +\infty$, $f(x) = \frac{x}{\eu^x} \to 0$
(\cref{thm:g12:exp:growth}); as $x \to -\infty$, $x \to -\infty$ and
$\eu^{-x} \to +\infty$, so $f(x) \to -\infty$.

$f''(x) = -\eu^{-x} - (1-x)\eu^{-x} = (x - 2)\eu^{-x}$, which changes sign at
$x = 2$: \omterm{def:g12:deriv:inflection}{inflection point} at $\bigl(2,\, 2\eu^{-2}\bigr)$.
\end{solution}

\begin{solution}{exo:g12:exp:5}
Let $g(x) = x - \ln(1+x)$ on $\intoo{-1}{+\infty}$. Then
$g'(x) = 1 - \frac{1}{1+x} = \frac{x}{1+x}$, negative on $\intoo{-1}{0}$ and
positive on $\intoo{0}{+\infty}$: $g$ attains its \omterm{def:g10:functions:extrema}{minimum} $g(0) = 0$, so
$g \geq 0$ with equality only at $0$. Hence $\ln(1+x) \leq x$.

Apply this with $x = \frac1n$:
$n\ln\left(1 + \frac1n\right) \leq 1$, so
$\left(1+\frac1n\right)^n = \eu^{\,n\ln(1+1/n)} \leq \eu$.
Apply it with $x = -\frac{1}{n+1} > -1$:
$\ln\left(1 - \frac{1}{n+1}\right) \leq -\frac{1}{n+1}$, so
$-(n+1)\ln\left(1 - \frac{1}{n+1}\right) \geq 1$ and
$\left(1 - \frac{1}{n+1}\right)^{-(n+1)} \geq \eu$.
\end{solution}

\begin{solution}{exo:g12:exp:6}
\emph{1.} Each year multiplies the capital by $1.03$: $C_n = C_0 (1.03)^n$.

\emph{2.} $C_n \geq 2C_0 \iff (1.03)^n \geq 2 \iff n \ln 1.03 \geq \ln 2
\iff n \geq \frac{\ln 2}{\ln 1.03} \approx 23.45$: the capital doubles after
$24$ years.

\emph{3.} $\frac{70}{3} \approx 23.3$, close to the exact
$\frac{\ln 2}{\ln 1.03}$. The rule works because
$\ln(1 + r) \approx r$ for small $r$, so
$\frac{\ln 2}{\ln(1+r)} \approx \frac{0.693}{r} \approx \frac{70}{100r}$.
\end{solution}

\begin{solution}{exo:g12:exp:7}
$\eu^{2x} > \eu^{x+1} \iff 2x > x + 1 \iff x > 1$ (the exponential is
strictly \omterm{def:g12:seq:monotonic}{increasing}): $S = \intoo{1}{+\infty}$.

\omterm{def:g11:func:function}{Domain} of the second inequality: $x^2 - 1 > 0$ and $x + 5 > 0$, \emph{i.e.}
$x \in \intoo{-5}{-1} \cup \intoo{1}{+\infty}$. On this \omterm{def:g11:func:function}{domain}, $\ln$ being
strictly \omterm{def:g12:seq:monotonic}{increasing},
\[
\ln(x^2-1) \leq \ln(x+5) \iff x^2 - 1 \leq x + 5 \iff x^2 - x - 6 \leq 0
\iff x \in \intcc{-2}{3}.
\]
Intersecting with the \omterm{def:g11:func:function}{domain}: $S = \intco{-2}{-1} \cup \intoc{1}{3}$.
\end{solution}

\begin{solution}{exo:g12:exp:8}
\emph{1.} $f'(x) = \dfrac{\eu^x x - \eu^x}{x^2} = \dfrac{(x-1)\eu^x}{x^2}$,
negative on $\intoo{0}{1}$, positive on $\intoo{1}{+\infty}$: \omterm{def:g10:functions:extrema}{minimum}
$f(1) = \eu$. Limits: $f \to +\infty$ at $0^+$ (numerator $\to 1$,
denominator $\to 0^+$) and at $+\infty$ (\cref{thm:g12:exp:growth}).

\emph{2.} For $x > 0$, $\eu^x = kx \iff f(x) = k$. From the \omterm{def:g10:functions:table}{variation table}
($+\infty \searrow \eu \nearrow +\infty$, \omterm{def:g12:limcont:continuity}{continuous} and strictly \omterm{def:g12:seq:monotonic}{monotonic}
on each side): no solution for $k < \eu$; exactly one ($x = 1$) for
$k = \eu$; exactly two for $k > \eu$ (one in $\intoo{0}{1}$, one in
$\intoo{1}{+\infty}$, by the bijection theorem on each \omterm{def:g10:numbers:interval}{interval}).
\end{solution}

\begin{solution}{exo:g12:exp:9}
\emph{1.} Direct from the first inequality of \cref{exo:g12:exp:5}:
$u_n = \eu^{\,n\ln(1+1/n)} \leq \eu^1 = \eu$.

\emph{2.} Write
\[
\ln u_n = n \ln\!\left(1 + \frac1n\right)
= \frac{\ln\!\left(1 + \frac1n\right) - \ln 1}{\frac1n},
\]
a difference quotient of $\ln$ at the point $1$ with increment
$h = \frac1n \to 0$. Since $\ln$ is \omterm{def:g12:deriv:derivative}{differentiable} at $1$ with \omterm{def:g12:deriv:derivative}{derivative}
$1$, $\ln u_n \to 1$. By \omterm{def:g12:limcont:continuity}{continuity} of $\exp$
(\cref{prop:g12:limcont:seqcont}), $u_n = \eu^{\ln u_n} \to \eu^1 = \eu$.
\end{solution}

\begin{solution}{pb:g12:exp:1}
\textbf{1.} $x = \frac{\ln 5}{2}$. \quad
$3x - 1 = \eu^2$: $x = \frac{\eu^2 + 1}{3}$. \quad With
$u = \eu^x$: $u^2 - 3u + 2 = (u - 1)(u - 2) = 0$: $x = 0$ or
$x = \ln 2$.

\textbf{2.} $\frac{3\ln 2}{\ln 2} = 3$; \quad
$3 + \frac12 = \frac72$; \quad $\frac52$.

\textbf{3.} $f'(x) = \eu^x - 1$: negative before $0$, positive
after: \omterm{def:g10:functions:extrema}{minimum} $f(0) = 0$, so $f \geq 0$ everywhere:
$\eu^x \geq 1 + x$. Substituting $x = \ln(1 + u)$:
$1 + u \geq 1 + \ln(1 + u)$, i.e.\ $\ln(1 + u) \leq u$.

\textbf{4.} $x^2\eu^{-x} = \frac{x^2}{\eu^x} \to 0$ and
$\frac{\ln x}{x} \to 0$: exponentials crush powers, powers
crush \omterm{def:g12:exp:ln}{logarithms}. And $\frac{\eu^x - 1}{x} =
\frac{\eu^x - \eu^0}{x - 0} \to \exp'(0) = 1$.

\textbf{5.} $f'(x) = \ln x + 1$: zero at $\eu^{-1}$, negative
before, positive after: \omterm{def:g10:functions:extrema}{minimum}
$f\!\left(\frac1\eu\right) = -\frac1\eu$; and
$x \ln x \to 0$ at $0^+$: the curve leaves the origin, dips to
$-\frac1\eu$, and climbs away.

\textbf{6.} $n = \frac{\ln 2}{\ln 1.03} \approx 23.4$ years.

\textbf{7.} Exact: $23.4$, $11.9$, $8.0$ years; rule of 72:
$24$, $12$, $8$. Since $\ln(1+r) \approx r$,
$\frac{\ln 2}{\ln(1+r)} \approx \frac{0.693}{r}$, i.e.\
$\frac{69.3}{\text{rate in }\%}$; bankers round up to $72$
because it divides beautifully by $2, 3, 4, 6, 8, 9, 12$ ---
mental arithmetic beats a decimal of accuracy.

\textbf{8.} $2^{t/20} = \eu^{t \ln 2 / 20}$:
$\lambda = \frac{\ln 2}{20}$ per minute. After
$24 \times 60 = 1440$ minutes: $2^{72} \approx 4.7 \times
10^{21}$ bacteria --- more than the grains of sand on Earth,
from one cell in one day. The model is honest only while food
and space last: real growth bends into the logistic S-curve
(the epidemiologist's curve of \cref{pb:g12:deriv:1}).

\textbf{9.} At 23:00 ($8$ hours):
$100 \times 2^{-8/5} \approx 33$ mg. Below $10$ mg:
$2^{-t/5} < 0.1$ gives $t > 5\,\frac{\ln 10}{\ln 2} \approx
16.6$ h: around 7:40 the next morning --- the espresso at
three has a long tail.

\textbf{10.} Yearly: $2$. Monthly:
$\left(1 + \frac{1}{12}\right)^{12} \approx 2.613$. Daily:
$\approx 2.715$. The ceiling: $\eu = 2.71828\ldots$
(\cref{exo:g12:exp:9}) --- compounding \omterm{def:g12:limcont:continuity}{continuously}, greed
\omterm{def:g12:seq:limit}{converges}.

\textbf{11.} If cases grow exponentially, $c(t) = c_0
\eu^{\lambda t}$, then $\ln c(t) = \ln c_0 + \lambda t$: a
straight line of \omterm{def:g10:reffunc:affine}{slope} $\lambda$ --- the growth rate. A
straight stretch on the log plot \emph{is} the exponential
phase, and its steepness is the epidemic's tempo.

\textbf{12.} $120 - 60 = 60$ dB means
$10\log_{10}(I_2/I_1) = 60$: the concert is
$10^6$ --- a million times --- more intense than the
conversation.

\textbf{13.} Amplitude: $10^{7-5} = 100$. Energy:
$100^{3/2} = 1000$: two magnitude points hide three orders of
magnitude in energy.

\textbf{14.} $10^{6.5 - 2} = 10^{4.5} \approx 32\,000$: lemon
juice is thirty thousand times more acidic than milk --- pH
compresses chemistry's chasms into a pocket scale.

\textbf{15.} $\log_2\!\frac{4186}{27.5} =
\frac{\ln(4186/27.5)}{\ln 2} \approx 7.25$: a piano spans just
over seven octaves. Senses respond to \emph{ratios} of stimuli
--- doubling the intensity feels like one step, whatever the
starting level --- so perception is built on a logarithmic
scale, and so are the units we invented for it.

\textbf{16.} Placing the stick for $a$ end-to-end with the
stick position for $b$ adds the lengths $\ln a + \ln b$, and
the graduation sitting at that total length reads
$\eu^{\ln a + \ln b} = ab$: the slide rule computes products
by adding \omterm{def:g12:exp:ln}{logarithms} --- $\ln(ab) = \ln a + \ln b$
(\cref{prop:g12:exp:lnalg}) carved in boxwood.

\textbf{17.} The \omterm{def:g10:functions:extrema}{minimum} of $\frac{\eu^x}{x}$ on
$\intoo{0}{+\infty}$ is $\eu$, at $x = 1$: no solution for
$k < \eu$, exactly one for $k = \eu$, two for $k > \eu$.
\omterm{def:g12:deriv:tangent}{Tangent} check: at $a = 1$ the \omterm{def:g12:deriv:tangent}{tangent} to $\eu^x$ is
$y = \eu + \eu(x - 1) = \eu x$: it passes through the origin
--- the threshold line, grazing the curve at $(1, \eu)$.

\textbf{18.} $n \ln\!\left(1 - \frac1n\right) =
\frac{\ln(1 - 1/n)}{-1/n} \times (-1) \to -1$, so
$\left(1 - \frac1n\right)^n \to \eu^{-1} \approx 0.368$: the
$0.999^{1000} \approx 0.368$ of \cref{pb:g11:binom:1},
explained --- the lottery's near-miss constant is
$\frac1\eu$.

\textbf{19.} Both are the limit shape of ``many independent
\omterm{def:g11:prob:model}{events}, each individually unlikely'': the chance that
\emph{none} of $n$ chances of size $\frac1n$ fires tends to
$\eu^{-1}$, and the secretary rule tunes its rejection window
so that success concentrates at that same constant.
$\frac1\eu \approx 0.368$.

\textbf{20.} The ceiling of compounding; the unique base whose
\omterm{def:g12:deriv:tangent}{tangent} at $0$ has \omterm{def:g10:reffunc:affine}{slope} $1$ (which is why calculus loves it);
the growth that outruns every power; the constant where
near-misses and optimal stopping settle; and the \omterm{def:g12:exp:ln}{logarithm} ---
multiplication become addition, the world's wildest ranges
folded onto a ruler.
\end{solution}
